% Options for packages loaded elsewhere
\PassOptionsToPackage{unicode}{hyperref}
\PassOptionsToPackage{hyphens}{url}
\documentclass[
]{article}
\usepackage{xcolor}
\usepackage{graphicx}
\usepackage{amsmath,amssymb}
\setcounter{secnumdepth}{-\maxdimen} % remove section numbering
\usepackage{iftex}
\ifPDFTeX
  \usepackage[T1]{fontenc}
  \usepackage[utf8]{inputenc}
  \usepackage{textcomp} % provide euro and other symbols
\else % if luatex or xetex
  \usepackage{unicode-math} % this also loads fontspec
  \defaultfontfeatures{Scale=MatchLowercase}
  \defaultfontfeatures[\rmfamily]{Ligatures=TeX,Scale=1}
\fi
\usepackage{lmodern}
\ifPDFTeX\else
  % xetex/luatex font selection
\fi
% Use upquote if available, for straight quotes in verbatim environments
\IfFileExists{upquote.sty}{\usepackage{upquote}}{}
\IfFileExists{microtype.sty}{% use microtype if available
  \usepackage[]{microtype}
  \UseMicrotypeSet[protrusion]{basicmath} % disable protrusion for tt fonts
}{}
\makeatletter
\@ifundefined{KOMAClassName}{% if non-KOMA class
  \IfFileExists{parskip.sty}{%
    \usepackage{parskip}
  }{% else
    \setlength{\parindent}{0pt}
    \setlength{\parskip}{6pt plus 2pt minus 1pt}}
}{% if KOMA class
  \KOMAoptions{parskip=half}}
\makeatother
\usepackage{longtable,booktabs,array}
\usepackage{caption}
\captionsetup[table]{skip=6pt}
\newcounter{none} % for unnumbered tables
\usepackage{calc} % for calculating minipage widths
% Correct order of tables after \paragraph or \subparagraph
\usepackage{etoolbox}
\makeatletter
\patchcmd\longtable{\par}{\if@noskipsec\mbox{}\fi\par}{}{}
\makeatother
\setlength{\emergencystretch}{3em} % prevent overfull lines
\providecommand{\tightlist}{%
  \setlength{\itemsep}{0pt}\setlength{\parskip}{0pt}}
\usepackage{bookmark}
\IfFileExists{xurl.sty}{\usepackage{xurl}}{} % add URL line breaks if available
\urlstyle{same}
% fallback for those not using the hyperref driver hyperxmp:
\makeatletter
\@ifundefined{xmpquote}{\newcommand{\xmpquote}[1]{#1}}{}
\makeatother
\hypersetup{
  hidelinks,
  pdfcreator={LaTeX via pandoc}}

\author{}
\date{}

\begin{document}

\section{Making Critical DNA Evidence Accessible: What Current
Statistics Can Leave Unsaid for the
Defense}\label{making-critical-dna-evidence-accessible-what-current-statistics-can-leave-unsaid-for-the-defense}

\subsection{With an AI Assistant
rob.v1}\label{with-an-ai-assistant-rob.v1}

Monte Miller, PhD; Forensic DNA Experts LLC

\subsection{Abstract}\label{abstract}

A correctly calculated DNA statistic need not answer the proposition a
factfinder is being asked to decide. This applied-mathematical article
connects three questions: which contributor configurations a comparison
evaluates, what the observations resolve about genetic states, and which
material conclusions a report expressly communicates. In a strict
two-contributor construction, two named people receive separate
likelihood ratios exceeding \(10^{13}\) against a two-unknown
alternative, while their complete-pair likelihood ratio is zero. The
comparisons address different configurations; they are not exhaustive
assessments of individual contribution. The same construction shows that
large relative support can coexist with diffuse genotype uncertainty.
Further synthetic examples distinguish population rarity, compatibility,
and quantitative discrimination. A comparison of three reports, holding
disclosed facts constant, separates derivability, a warning about
evaluative scope, and an express joint assessment. Building on
established assessment and reporting literature, the article identifies
different responses to different overinterpretations: another
proposition-specific assessment, clarification of genotype uncertainty,
explicit communication of an available conclusion, or reevaluation under
another model. It proposes a bounded transparency workflow rather than
another statistic for every report. A short legal discussion advances,
without attributing a holding, the unresolved hypothesis that McDaniel
v. Brown's concern with unsupported probabilistic meaning may also bear
on overstatement of a comparison's scope. Exact inputs and an
accompanying implementation reproduce the central construction and six
displayed quantitative settings. Operational validity, improved
comprehension, reporting prevalence, and case-specific unfairness are
not established.

\textbf{Keywords:} forensic DNA; likelihood ratio; joint contribution;
genotype uncertainty; proposition selection; evidence communication;
population conditioning; DNA reporting; DNA LR prior odds; DNA LR
posterior odds.

\section{1. Introduction}\label{introduction}

The study began by asking one question: Is it reasonable that every
legal DNA sample is accurately and adequately conveyed by a single
statistic? The forensic DNA likelihood ratio from probabilistic
genotyping is typically the only statistic given and discussed once the
lab report is produced, and while it might be helpful, this manuscript
endorses providing the electropherogram (EP) along with additional known
statistics that specifically describe it for further reference. This
manuscript takes the perspective that someone must wager that the DNA
from the person of interest (POI) is, in fact, on an item of evidence,
but they can only bet if convinced beyond a reasonable doubt about
winning that wager, because they wager the future of the POI's life. The
perspective pursued is one where the scientists are outside of the
deliberation, and are attempting to provide whatever forensic DNA
information is available that will enable the bettor to make as fully an
informed decision as possible. The submitted manuscript is an applied
mathematics study, and it presents nothing mathematically original or
novel. It is built with the statistics of others by modularly stacking
their mathematics. This is done in order to help directly explain a DNA
EP under a variety of conditions; how to mathematically address the
conditions presented is not new. It is not any specific statistic that
is compelling in the different scenarios herein, instead, a more
informed decision is made possible by the synergy and context of
examining what is germane and salient under different conditions, and
presenting that information as a whole, to explain the EP and DNA
evidence to the bettor. A DNA statistic becomes meaningful through the
question it answers. A population frequency concerns a genetic event
under a population model. A genotype posterior concerns uncertainty over
genetic states under observations and a prior. A likelihood ratio
compares observations under stated propositions. None automatically
establishes personal contribution, activity, or guilt.

The difficulty is not always incorrect arithmetic. A calculation can be
correct while an inference drawn from it concerns a different
proposition. A report can also provide the facts needed to derive a
material limitation without expressly reporting the limitation or
whether the relevant question was evaluated.

Consider the construction developed below. At each of thirteen
independently stipulated loci, the observed allele set is
\(\{1,2,3,4\}\). Person A carries \(2/3\) and person B carries \(1/3\).
A fits with an unknown partner carrying \(1/4\); B fits with an unknown
carrying \(2/4\). Together, however, A and B lack allele \(4\).

Under the stated population and observation models, their separate
comparisons yield approximately \(6.69\times10^{13}\) and
\(3.80\times10^{14}\) against two unknown contributors. The complete
named pair has likelihood ratio zero. The results do not contradict one
another: their numerators describe different configurations.

\subsection{1.1 Audience, precedents, and
contribution}\label{audience-precedents-and-contribution}

This article addresses forensic scientists, statistical researchers, and
legal professionals who evaluate or communicate DNA evidence. It is an
academic methodological analysis, not a validated courtroom presentation
or a standalone reporting template.

Previous work establishes the importance of evaluating multiple persons
of interest jointly and addresses how the resulting assessments should
be reported. Slooten (2022) develops a hypothesis-based approach to
evaluation and reporting, emphasizing the likelihoods of the relevant
hypotheses and the additional weights needed for composite
individual-contribution comparisons. Duke et al.~(2022) investigate
compound and conditioned likelihood-ratio behavior. Wivell et al.~(2023)
expressly recommend reporting when candidates cannot jointly be donors.
The ASB 041 ballot draft also addresses proposition selection,
attribution to legal parties, joint assessment, and reassessment.

The overlap therefore concerns reporting as well as calculation.

{\def\LTcaptype{none} % do not increment counter
\begin{longtable}[]{@{}
  >{\raggedright\arraybackslash}p{(\linewidth - 4\tabcolsep) * \real{0.3333}}
  >{\raggedright\arraybackslash}p{(\linewidth - 4\tabcolsep) * \real{0.3333}}
  >{\raggedright\arraybackslash}p{(\linewidth - 4\tabcolsep) * \real{0.3333}}@{}}
\toprule\noalign{}
\begin{minipage}[b]{\linewidth}\raggedright
Precedent
\end{minipage} & \begin{minipage}[b]{\linewidth}\raggedright
Relevant contribution
\end{minipage} & \begin{minipage}[b]{\linewidth}\raggedright
Relationship to this article
\end{minipage} \\
\midrule\noalign{}
\endhead
\bottomrule\noalign{}
\endlastfoot
Slooten (2022), Introduction & Organizes evaluation around relevant
hypotheses and their likelihoods; discusses reporting likelihood tables
and the weights needed for composite individual-contribution LRs. & The
proposition map below is a compact use of an established approach, not a
new attribution method. \\
Duke et al.~(2022), Introduction and Sections 3.3--3.7 & Investigates
compound and conditioned LR behavior, including sampling-related zeros
and atypical peak effects. & The exact structural zero below must be
distinguished from computational or observation-model problems. \\
Wivell et al.~(2023), Introduction and Section 5 & Discusses
alternatives when the defense offers no proposition, joint versus
individual evidential weight, and explicit reporting of joint
incompatibility. & Explicit joint reporting and caution about party
labels are acknowledged precedents. \\
AAFS Standards Board (2021), ASB 041 ballot draft, Foreword and clauses
4.3--4.5, 4.14--4.16 & Addresses party attribution, limitations of
chosen propositions, joint assessment, additional comparisons, and
reasonable reassessment requests. & The workflow proposed here shares
this ground. The cited document is a ballot draft, not evidence of a
binding or universally adopted requirement. \\
\end{longtable}
}

The present contribution is a linked applied-mathematical explanation of
three questions:

\begin{enumerate}
\def\labelenumi{\arabic{enumi}.}
\tightlist
\item
  Which contributor configurations does a comparison evaluate?
\item
  What do the observations resolve about genetic states?
\item
  Which material conclusions does a report expressly communicate?
\end{enumerate}

The strict construction links two distinct findings: strong support in
selected comparisons can coexist with complete-pair incompatibility, and
large relative support can coexist with diffuse genotype uncertainty.
Example D changes signal strength while holding target rarity fixed; G
examines changing quantitative patterns with population-weighted unknown
co-contributors. The reporting comparison then holds disclosed facts
constant to distinguish a derivable conclusion, a warning about
evaluative scope, and an expressly reported assessment.

The analytical purpose is to identify which response a particular
overinterpretation requires. A joint-contribution claim may require
another proposition-specific assessment. A claim of genotype resolution
may require examination of genotype uncertainty. An unstated material
conclusion may require explicit reporting rather than another
calculation. A conclusion transferred to a different observation
mechanism requires reevaluation under that mechanism.

These responses are not interchangeable, and adding another statistic is
not a universal remedy. The article develops this integration without
claiming a new joint-assessment method, a historical first for each
distinction, or an absence of earlier reporting recommendations.

\subsection{1.2 Hierarchy of
conclusions}\label{hierarchy-of-conclusions}

{\def\LTcaptype{none} % do not increment counter
\begin{longtable}[]{@{}
  >{\raggedright\arraybackslash}p{(\linewidth - 4\tabcolsep) * \real{0.3333}}
  >{\raggedright\arraybackslash}p{(\linewidth - 4\tabcolsep) * \real{0.3333}}
  >{\raggedright\arraybackslash}p{(\linewidth - 4\tabcolsep) * \real{0.3333}}@{}}
\toprule\noalign{}
\begin{minipage}[b]{\linewidth}\raggedright
Evidentiary status
\end{minipage} & \begin{minipage}[b]{\linewidth}\raggedright
Component
\end{minipage} & \begin{minipage}[b]{\linewidth}\raggedright
Conclusion and boundary
\end{minipage} \\
\midrule\noalign{}
\endhead
\bottomrule\noalign{}
\endlastfoot
Demonstrated mathematical results & Strict construction; matched
posterior--LR identities; population-filter calculations & Results
follow under the specified models. Application to an actual case
requires separate justification. \\
Explanatory illustrations & Synthetic quantitative examples D and G &
They distinguish rarity, discrimination, and genotype uncertainty. They
do not validate laboratory methods or establish casework prevalence. \\
Constructed reporting analysis & Versions A--C under common disclosed
facts & Derivability, a scope warning, and an express joint assessment
are different report features. Reader effects are not measured. \\
Normative recommendation & Explicit communication of material
assessments and limitations & The justification is transparency, not a
demonstrated comprehension benefit or a universal legal obligation. \\
Unresolved applications and interpretation & Implementation-specific
conditioning, actual reporting practices, and the proposed McDaniel
connection & These require additional evidence or argument. They are not
additional demonstrated defects. \\
\end{longtable}
}

The organizing question is whether the assessment calculated supports
the meaning attributed to it. A six-face analogy---observations, model,
propositions, mathematics, reporting, and legal context---emphasizes
their connections. As with a partially aligned Rubik's cube, local
correctness does not establish coherence of the whole arrangement. The
analogy organizes the inquiry; it does not prove alignment or determine
legal fairness.

A definite result under stated assumptions can coexist with uncertainty
about those assumptions' suitability for an application. Conversely,
questioning applicability does not establish that a particular
implementation or report is defective.

\begin{center}\rule{0.5\linewidth}{0.5pt}\end{center}

\section{Part I --- Propositions and mathematical
construction}\label{part-i-propositions-and-mathematical-construction}

\section{2. Essential definitions and
conditioning}\label{essential-definitions-and-conditioning}

Let \(E\) denote the questioned observations and \(I\) the relevant
background information. A genotype is an unordered diploid allele pair
at one locus. A multilocus profile is a vector of such genotypes.
Different people can share a genotype or profile; a genetic state is not
a uniquely identified person.

A likelihood ratio compares evidence likelihoods:

\[
\operatorname{LR}(H_1:H_0)
=
\frac{\Pr(E\mid H_1,I)}
{\Pr(E\mid H_0,I)}.
\]

Densities replace point probabilities for continuous observations. A
finite reported ratio requires a positive denominator.

An LR does not establish which proposition is true. Posterior odds
require prior odds. Contribution assumed within a proposition remains
conditional; it is not an established identity. Whether a legal burden
is satisfied is a further question.

\subsection{2.1 A common genotype-replacement
comparison}\label{a-common-genotype-replacement-comparison}

Let \(\mathcal G\) be the declared genetic state space, \(t\) a target
state, and \(\pi(g)\) the distribution assigned to an alternative
person's genetic state under the stated population and background
assumptions.

Unless expressly identified as multilocus, states are single-locus
genotypes. For a designated contributor position and contributor count,
let \(L(g)\) be the probability or density of \(E\) when that position's
genotype is fixed at \(g\), conditional on \(I\). It includes the
declared treatment of co-contributor genotypes and other nuisance
quantities.

A common-model comparison fixes \(t\) in that position and replaces it
under the alternative with a population draw:

\[
R=
\frac{L(t)}
{\sum_{g\in\mathcal G}\pi(g)L(g)}.
\]

In these examples,

\[
f(t)=\pi(t).
\]

A person-level interpretation requires person propositions that justify
this construction. Fixing a genotype in a position, stipulating that a
named person occupies that position, and assessing whether the person
contributed anywhere are distinct statements.

Identities using a common \(L(g)\) require compatible state spaces,
conditioning, and co-contributor and nuisance treatment.

\subsection{2.2 Population assumptions}\label{population-assumptions}

Under the stipulated Hardy--Weinberg model,

\[
f(a/b)=
\begin{cases}
2p_ap_b,&a\ne b,\\
p_a^2,&a=b.
\end{cases}
\]

For a multilocus profile, a product representation requires appropriate
factorization:

\[
f(t)=\prod_{\ell=1}^{m}f_\ell(t_\ell).
\]

Allele frequencies, Hardy--Weinberg proportions, between-person
independence, population relevance, and temporal stability are separate
assumptions. Hardy--Weinberg proportions alone do not establish
independence between relatives or justify using one population
distribution for every alternative contributor. Population structure and
relationships can require different conditional distributions (National
Research Council, 1996, Chapter 4). Reference-selection circumstances
also require separate consideration, as discussed in Section 2.3.

The examples stipulate frequencies; they do not estimate temporal change
or local-population composition.

\subsection{2.3 Different operations called
conditioning}\label{different-operations-called-conditioning}

{\def\LTcaptype{none} % do not increment counter
\begin{longtable}[]{@{}
  >{\raggedright\arraybackslash}p{(\linewidth - 4\tabcolsep) * \real{0.3333}}
  >{\raggedright\arraybackslash}p{(\linewidth - 4\tabcolsep) * \real{0.3333}}
  >{\raggedright\arraybackslash}p{(\linewidth - 4\tabcolsep) * \real{0.3333}}@{}}
\toprule\noalign{}
\begin{minipage}[b]{\linewidth}\raggedright
Operation
\end{minipage} & \begin{minipage}[b]{\linewidth}\raggedright
What is specified
\end{minipage} & \begin{minipage}[b]{\linewidth}\raggedright
Consequence
\end{minipage} \\
\midrule\noalign{}
\endhead
\bottomrule\noalign{}
\endlastfoot
Treating a reference genotype as background & An observed genetic state
is known. & Under fixed frequencies and independence, it need not change
another unrelated person's genotype distribution. \\
Assuming a person contributed & A proposition assigns that person to the
contributor configuration. & The evidence likelihood changes through
that configuration, even if population frequencies remain fixed. \\
Conditioning on a family relationship & A pedigree or relationship
connects observed and unobserved genotypes. & Unknown-genotype
predictions can depend on observed family genotypes. \\
Updating shared population uncertainty & Relevant people share uncertain
population parameters within a joint model. & Reference observations can
affect predictive genotype distributions if their selection and
relevance are modeled appropriately. \\
\end{longtable}
}

Observed genotypes are not altered by these operations; distributions
over unknown states may change.

Additional references are not automatically unbiased population samples.
Relatives, randomly sampled residents, and people selected through a DNA
match provide different information. Geographic proximity or the tested
person's ancestry does not, by itself, determine the appropriate
alternative-contributor distribution.

Questioned-mixture alleles are evidence to explain, not automatic
independent population draws. Reference information and selection
circumstances must be treated coherently across propositions without
double counting.

When LRs differ, identify whether the difference arises from
propositions, observations, background, population or relationship
distributions, or nuisance treatment. Numerical disagreement does not
establish an error, and numerical agreement does not establish identical
meaning.

This conditioning framework is conceptual. It establishes no particular
correction, direction of LR change, or implementation defect.

\section{3. Strong separate support with complete-pair
incompatibility}\label{strong-separate-support-with-complete-pair-incompatibility}

\subsection{3.1 Strict construction}\label{strict-construction}

Consider exactly two diploid contributors and an observation rule
recording their complete allele set without dropout, drop-in, or error.

At every stipulated locus,

\[
E_\ell=\{1,2,3,4\}.
\]

A has genotype \(2/3\) and B has genotype \(1/3\). Frequencies for
alleles \(1\) through \(5\) are

\[
(0.07,0.08,0.09,0.10,0.66).
\]

With four distinct observed alleles and two diploid contributors, every
fitting configuration partitions the four alleles into two disjoint
heterozygotes. There are three unordered partitions and six ordered
assignments.

A requires a partner carrying \(1/4\). B requires a partner carrying
\(2/4\). Together A and B supply alleles \(1\), \(2\), and \(3\), but
not allele \(4\).

\subsection{3.2 Proposition map}\label{proposition-map}

Under this two-person model, the following scenarios distinguish the
named-person inclusion possibilities:

{\def\LTcaptype{none} % do not increment counter
\begin{longtable}[]{@{}
  >{\raggedright\arraybackslash}p{(\linewidth - 4\tabcolsep) * \real{0.3333}}
  >{\raggedright\arraybackslash}p{(\linewidth - 4\tabcolsep) * \real{0.3333}}
  >{\raggedright\arraybackslash}p{(\linewidth - 4\tabcolsep) * \real{0.3333}}@{}}
\toprule\noalign{}
\begin{minipage}[b]{\linewidth}\raggedright
Scenario
\end{minipage} & \begin{minipage}[b]{\linewidth}\raggedright
Complete contributor configuration
\end{minipage} & \begin{minipage}[b]{\linewidth}\raggedright
Unknown-person convention
\end{minipage} \\
\midrule\noalign{}
\endhead
\bottomrule\noalign{}
\endlastfoot
\(H_{AB}\) & A and B & Both named people form the complete pair. \\
\(H_{AU}\) & A and one other person & The other person is neither A nor
B. \\
\(H_{BU}\) & B and one other person & The other person is neither A nor
B. \\
\(H_{UU}\) & Two other people & They are distinct people, neither of
whom is A or B. \\
\end{longtable}
}

Identity exclusion does not imply genotype exclusion. Under the
fixed-frequency independent-person model, unknown genotypes remain
population-distributed and may match either reference. Both references
are background information under every scenario.

Writing

\[
L_{XY}=\Pr(E\mid H_{XY},I),
\]

the displayed comparisons are

\[
R_A=\frac{L_{AU}}{L_{UU}},\qquad
R_B=\frac{L_{BU}}{L_{UU}},\qquad
R_{AB}=\frac{L_{AB}}{L_{UU}}.
\]

The first two are selected comparisons. A composite assessment of ``A
contributed'' would combine \(H_{AB}\) and \(H_{AU}\); its
non-contribution alternative would combine \(H_{BU}\) and \(H_{UU}\).
Calculating that comparison requires specified distributions within each
side. The separate LR \(R_A\) does not supply them. The corresponding
distinction applies to B.

This follows the established distinction between likelihoods for
specified configurations and composite individual-contribution
comparisons discussed by Slooten (2022). No equal weighting of the four
scenarios is assumed here.

\subsection{3.3 Positions and
assignment}\label{positions-and-assignment}

The contributor positions are exchangeable bookkeeping labels, not
empirically resolved major and minor roles. Let \(L_{A,1}\) and
\(L_{A,2}\) be the full-profile likelihoods with A assigned to the first
and second positions, respectively, and a population-weighted unknown in
the other position. Symmetry gives

\[
L_{A,1}=L_{A,2}.
\]

Prespecified assignment weights therefore give

\[
\lambda L_{A,1}+(1-\lambda)L_{A,2}=L_{A,1},
\qquad 0\leq\lambda\leq1.
\]

This is a mixture over assignments, not the sum of two unweighted copies
of the same likelihood. The assignment concerns the whole profile, not a
different named-person position selected independently at each locus.

The equality is model-specific. It does not provide a general conversion
for fraction-linked positions.

\subsection{3.4 Separate and joint
results}\label{separate-and-joint-results}

Each compatible ordered assignment has probability \(4p_1p_2p_3p_4\)
under two independent population draws. Thus the one-locus denominator
is

\[
D=24p_1p_2p_3p_4=0.0012096.
\]

The separate numerator probabilities are

\[
L_A=2p_1p_4=0.014,
\qquad
L_B=2p_2p_4=0.016.
\]

Here \(D\), \(L_A\), and \(L_B\) denote the one-locus likelihoods for
\(H_{UU}\), \(H_{AU}\), and \(H_{BU}\), respectively.

{\def\LTcaptype{none} % do not increment counter
\begin{longtable}[]{@{}
  >{\raggedright\arraybackslash}p{(\linewidth - 4\tabcolsep) * \real{0.2727}}
  >{\raggedleft\arraybackslash}p{(\linewidth - 4\tabcolsep) * \real{0.3636}}
  >{\raggedleft\arraybackslash}p{(\linewidth - 4\tabcolsep) * \real{0.3636}}@{}}
\toprule\noalign{}
\begin{minipage}[b]{\linewidth}\raggedright
Comparison against two unknowns
\end{minipage} & \begin{minipage}[b]{\linewidth}\raggedleft
One-locus LR
\end{minipage} & \begin{minipage}[b]{\linewidth}\raggedleft
Thirteen independent loci
\end{minipage} \\
\midrule\noalign{}
\endhead
\bottomrule\noalign{}
\endlastfoot
A with an unknown & Approximately \(11.5741\) & Approximately
\(6.69\times10^{13}\) \\
B with an unknown & Approximately \(13.2275\) & Approximately
\(3.80\times10^{14}\) \\
A and B as the complete pair & \(0\) & \(0\) \\
\end{longtable}
}

The complete pair fails the observation rule. The separate results
remain valid for their stated comparisons.

\begin{figure}
\centering
\includegraphics[width=\linewidth]{figure-1.png}
\caption{Separate compatible pairs and incompatible named pair}
\end{figure}

\textbf{Figure 1. Separate fit does not establish joint fit under the
strict two-person model.} A's genotype \(2/3\) requires partner \(1/4\),
whereas B's genotype \(1/3\) requires partner \(2/4\). Together, A and B
fail to supply observed allele \(4\). Each stated configuration is
compared with two unknown contributors under the specified
Hardy--Weinberg model. Both named references are background information;
unknown identities exclude A and B but not their genotypes. Positions
are exchangeable bookkeeping labels. Separate comparisons are selected
scenarios, not exhaustive individual-contribution assessments. Symbols
represent allele copies, not peak heights. Full-profile values refer to
thirteen independent loci, not merely the illustrated locus. The joint
zero concerns the complete pair under complete observation without
dropout, drop-in, or error; it is not exclusion under every larger
mixture or other model.

\subsection{3.5 What the incompatibility does and does not
establish}\label{what-the-incompatibility-does-and-does-not-establish}

The joint result concerns the complete named pair, not the
non-contribution of each individual. If the disputed proposition is that
A and B form the complete pair, it directly challenges that proposition.
If the dispute concerns A contributing with another person, it does not
resolve it.

A larger contributor configuration or another mechanism capable of
explaining the observed allele requires its own evaluation. Dropout
alone cannot supply allele \(4\), absent from both references.

A structural zero also differs from a computational failure to explore
supported configurations. Duke et al.~(2022) and Wivell et al.~(2023)
discuss sampling-related zeros in probabilistic-genotyping calculations;
Duke et al.~also examine atypical peak behavior. Their operational
examples do not determine the likelihoods in this strict model, and this
construction is not a contradiction of their empirical findings.

Thirteen loci amplify the contrast but are unnecessary for its
existence. If a repeatable experiment has selected-scenario LRs
\(r_A>1\) and \(r_B>1\), a zero joint likelihood, and a positive common
alternative likelihood, independent identically specified repetitions
give

\[
R_A^{(n)}=r_A^n,\qquad
R_B^{(n)}=r_B^n,\qquad
R_{AB}^{(n)}=0.
\]

This is a stipulated existence result, not a claim that arbitrarily many
real forensic loci are available or that marginal likelihoods remain
independent after integrating shared nuisance parameters.

\section{4. Large relative support and genotype
uncertainty}\label{large-relative-support-and-genotype-uncertainty}

\subsection{4.1 Matched posterior--likelihood
identities}\label{matched-posteriorlikelihood-identities}

For genotype prior \(q(g)\),

\[
w_q(g)=
\frac{q(g)L(g)}
{\sum_{h\in\mathcal G}q(h)L(h)},
\]

provided the normalizing likelihood is positive.

Under a uniform target-genotype prior,

\[
w_U(g)=\frac{L(g)}{\sum_{h\in\mathcal G}L(h)}.
\]

Under the alternative-person population prior,

\[
w_\pi(g)=
\frac{\pi(g)L(g)}
{\sum_{h\in\mathcal G}\pi(h)L(h)}.
\]

Consequently,

\[
R=\frac{w_\pi(t)}{f(t)},
\]

when \(f(t)>0\), the normalizing likelihood is positive, and state
spaces, conditioning, and conditional likelihoods are compatible.

This decomposes the same LR. It is not independent evidence, and
\(w_U(t)\) must not be substituted for \(w_\pi(t)\).

A contributor position is occupied under the assumed model. Its genotype
posterior describes uncertainty about that position's genetic state; it
does not establish which named person occupies it. The target frequency
is not generally the mixture evidence-likelihood denominator, which
averages observation likelihoods over states.

Recovery of likelihoods from exported weights requires additional
conditions, addressed in Appendix A.1. A weight-table heading alone does
not establish normalization, prior, completeness, state semantics, or
nuisance treatment.

\subsection{4.2 The strict construction}\label{the-strict-construction}

Under the two-unknown population model, the six compatible ordered
assignments have equal prior products. Each compatible genotype in a
designated position therefore has posterior probability \(1/6\) at one
locus.

For A,

\[
f_{A,\ell}=2p_2p_3,
\qquad
R_{A,\ell}=\frac{1}{12p_2p_3}
=\frac{1}{6f_{A,\ell}}.
\]

Similarly,

\[
R_{B,\ell}=\frac{1}{6f_{B,\ell}},
\qquad
f_{B,\ell}=2p_1p_3.
\]

Across thirteen independent loci,

\[
w_\pi(t)=6^{-13}\approx7.66\times10^{-11},
\]

for either specified compatible target profile in that position. There
are

\[
6^{13}=13{,}060{,}694{,}016
\]

compatible profiles for the position. These are genetic states, not
identifiable people.

Writing \(f_A\) and \(f_B\) for the target profile frequencies,

\[
R_A=\frac{6^{-13}}{f_A},
\qquad
R_B=\frac{6^{-13}}{f_B}.
\]

A small genotype posterior can be much larger than a still smaller
population probability. Large relative support and diffuse genotype
uncertainty are therefore compatible.

The posterior is relevant when the attributed claim is that the
observations resolve the target genetic state. Under the specified model
and prior, it describes how uncertainty remains distributed over genetic
states in the designated position. The LR instead measures relative
support for the stated comparison. These quantities are not rival
estimates of named-person contribution, and the matched identity
supplies neither a correction to the LR nor independent evidence.

Neither posterior is A's or B's contribution probability, and its
complement is not a probability of innocence.

\subsection{4.3 Different genotypes and different
people}\label{different-genotypes-and-different-people}

Another person may carry the target genotype. A comparison against other
genotypes therefore differs from a comparison against another person.

For \(0<f(t)<1\), let \(A_{\ne t}\) be the mean likelihood over
non-target genotypes under the alternative-person distribution
conditional on not carrying \(t\). The alternative-person likelihood is

\[
f(t)L(t)+\bigl(1-f(t)\bigr)A_{\ne t}.
\]

If \(L(t)>0\) and \(C=L(t)/A_{\ne t}\),

\[
R=
\frac{1}{f(t)+\bigl(1-f(t)\bigr)/C},
\]

with the limiting interpretation when \(A_{\ne t}=0\).

The matching-genotype branch is already part of the alternative, not a
missing independent factor. Since

\[
\sum_g\pi(g)L(g)\geq f(t)L(t),
\]

the matched comparison satisfies

\[
R\leq\frac{1}{f(t)}.
\]

This is not a universal bound for arbitrary person-level comparisons. It
requires the stated matching branch and common conditional likelihood.
If \(f(t)=1\), the common-model LR is \(1\) when the evidence likelihood
is positive. If \(f(t)=0\), use the underlying likelihoods directly; no
finite reciprocal-frequency bound follows.

For later illustrations, define the uniform-reference different-genotype
comparison

\[
B=
\frac{(M-1)L(t)}
{\sum_{g\ne t}L(g)},
\qquad M=|\mathcal G|>1.
\]

At fixed \(M\),

\[
B=\frac{(M-1)w_U(t)}{1-w_U(t)}.
\]

This is a transformation of the target weight for a specified
alternative, not additional independent evidence. Appendix A develops
its boundaries and explains why it does not generally determine the
different-person LR.

\begin{center}\rule{0.5\linewidth}{0.5pt}\end{center}

\section{Part II --- What the statistics count and
compare}\label{part-ii-what-the-statistics-count-and-compare}

\section{5. Population filters count different
sets}\label{population-filters-count-different-sets}

Return to the strict frequencies and observed set \(\{1,2,3,4\}\).

A compatibility rule defines a set \(\mathcal C(E)\) and assigns it
population mass

\[
P_{\mathrm{comp}}(E)=\sum_{g\in\mathcal C(E)}\pi(g).
\]

The rule defining the set is essential.

For target \(2/3\),

\[
f(2/3)=2p_2p_3=0.0144.
\]

Under the complete two-person rule, a contributor must carry one of the
six heterozygotes

\[
1/2,\quad1/3,\quad1/4,\quad2/3,\quad2/4,\quad3/4.
\]

Their total population probability is

\[
P_{\mathrm{comp},2}
=\sum_{1\leq i<j\leq4}2p_ip_j
=0.0862.
\]

Requiring only that both alleles lie in the observed set additionally
admits the four homozygotes. Its probability is

\[
P_{\mathrm{set}}
=(p_1+p_2+p_3+p_4)^2
=0.1156.
\]

A homozygote passes the outer rule but cannot, with one other diploid
contributor, explain four distinct observed alleles under the strict
model.

{\def\LTcaptype{none} % do not increment counter
\begin{longtable}[]{@{}llr@{}}
\toprule\noalign{}
Population event & States counted & Probability \\
\midrule\noalign{}
\endhead
\bottomrule\noalign{}
\endlastfoot
Exact target genotype & One genotype & \(0.0144\) \\
Strict two-person compatibility & Six heterozygotes & \(0.0862\) \\
Both alleles inside the observed set & Ten genotypes & \(0.1156\) \\
\end{longtable}
}

The squared allele-frequency sum is associated with single-locus
inclusion calculations in the combined-probability-of-inclusion
literature, subject to the assumptions governing eligible alleles and
loci (Bieber et al., 2016, Rule R5 and principle P1). The
six-heterozygote quantity here is explicitly called strict two-person
compatibility rather than assigned a guideline label.

\begin{figure}
\centering
\includegraphics[width=\linewidth]{figure-2.png}
\caption{Three nested population genotype filters}
\end{figure}

\textbf{Figure 2. Three population filters count different genotype
sets.} For observed alleles \(\{1,2,3,4\}\), target \(2/3\) has
population probability \(0.0144\). Under exactly two diploid
contributors, complete allele observation, and no dropout, drop-in, or
error, the six heterozygotes each admit a complementary partner; their
total population probability is \(0.0862\). Requiring only that both
alleles belong to the observed set also admits four homozygotes, giving
\(0.1156\). Probabilities use the stated Hardy--Weinberg model. Cell
areas do not represent probabilities, and outside-set genotypes are not
displayed. These are population events, not posterior probabilities of
contribution.

The differences arise from different counted events, not competing
estimates of one probability. A mechanism such as dropout may change the
relevant compatibility set; these rules must not be transferred without
reevaluation.

\section{6. Fixed rarity with changing quantitative signal: example
D}\label{fixed-rarity-with-changing-quantitative-signal-example-d}

Examples D and G use a Gaussian observation model, not the strict rule.

D uses six allele categories with frequencies

\[
(0.30,0.24,0.18,0.13,0.09,0.06),
\]

giving twenty-one unordered genotypes.

The known major genotype is \(1/2\). The target in the designated minor
position is \(5/6\), with

\[
f(t)=2(0.09)(0.06)=0.0108.
\]

The target proposition fixes \(5/6\) in that position. The alternative
replaces it with a population-distributed genotype. The known major and
prescribed fractions remain the same.

Observations are deterministic synthetic signals. Evaluation uses
independent Gaussian channel errors with \(\sigma=0.035\). Fractions are
inputs, not estimates. These are settings at one toy locus, not
independent loci. Section 14 gives exact signals and the likelihood
specification.

{\def\LTcaptype{none} % do not increment counter
\begin{longtable}[]{@{}
  >{\raggedright\arraybackslash}p{(\linewidth - 8\tabcolsep) * \real{0.1579}}
  >{\raggedleft\arraybackslash}p{(\linewidth - 8\tabcolsep) * \real{0.2105}}
  >{\raggedleft\arraybackslash}p{(\linewidth - 8\tabcolsep) * \real{0.2105}}
  >{\raggedleft\arraybackslash}p{(\linewidth - 8\tabcolsep) * \real{0.2105}}
  >{\raggedleft\arraybackslash}p{(\linewidth - 8\tabcolsep) * \real{0.2105}}@{}}
\toprule\noalign{}
\begin{minipage}[b]{\linewidth}\raggedright
Setting
\end{minipage} & \begin{minipage}[b]{\linewidth}\raggedleft
Exact minor fraction
\end{minipage} & \begin{minipage}[b]{\linewidth}\raggedleft
Signal at each target allele
\end{minipage} & \begin{minipage}[b]{\linewidth}\raggedleft
\(R\)
\end{minipage} & \begin{minipage}[b]{\linewidth}\raggedleft
\(w_U(t)\)
\end{minipage} \\
\midrule\noalign{}
\endhead
\bottomrule\noalign{}
\endlastfoot
D1 & \(0\) & \(0\) & \(1\) & \(0.047619\) \\
D6 & \(1/9\) & \(1/18\) & \(27.9507\) & \(0.541729\) \\
D10 & \(1/5\) & \(1/10\) & \(91.9453\) & \(0.997158\) \\
\end{longtable}
}

At zero minor fraction, varying the genotype in that position changes no
predicted signal. Thus,

\[
R=1,\qquad
w_U(t)=\frac{1}{21},\qquad
w_\pi(t)=f(t).
\]

This is an uninformative boundary, not dropout, exclusion, or evidence
of a detectable second contributor.

As signal increases, the observations distinguish the target genotype
more strongly under the model, while its population rarity remains
unchanged.

\begin{figure}
\centering
\includegraphics[width=\linewidth]{figure-3.png}
\caption{Fixed rarity with changing minor signal}
\end{figure}

\textbf{Figure 3. Fixed genotype rarity with changing quantitative
discrimination in example D.} Three settings at one synthetic locus use
known major genotype \(1/2\), target \(5/6\) in the minor position, and
prescribed fractions. Heights represent supplied deterministic channel
signals. The LR \(R\) compares the target genotype in that position with
a population-distributed replacement, retaining the same major genotype,
fractions, and Gaussian model. The target population probability remains
\(0.0108\). The weight \(w_U\) uses a uniform target-genotype prior and
is not a named-person contribution probability. Zero minor fraction is
an uninformative boundary, not dropout or exclusion. Evaluation uses
independent Gaussian channel errors with \(\sigma=0.035\). Widths and
shapes are schematic, not simulated laboratory electropherograms.
Statistics are rounded; settings are not independent loci.

The choice of prior remains consequential. At D6,

\[
w_U(t)\approx0.541729,
\qquad
w_\pi(t)\approx0.301868.
\]

These are different model-dependent quantities, not inconsistent
estimates of one posterior.

\section{7. The same alleles with different quantitative information:
example
G}\label{the-same-alleles-with-different-quantitative-information-example-g}

G uses the same population frequencies and Gaussian error scale.
Synthetic signals are generated from

\[
5/6,\qquad3/4,\qquad1/2,
\]

in that fraction order.

During evaluation, all contributor genotypes are unknown except where
the target proposition fixes \(5/6\) in the first position. Other
genotypes are population-weighted. Under the alternative, all three
positions are population-weighted.

Prescribed fractions remain the same under both propositions. Unequal
fractions distinguish positions, so the strict construction's assignment
symmetry must not be assumed.

{\def\LTcaptype{none} % do not increment counter
\begin{longtable}[]{@{}llrr@{}}
\toprule\noalign{}
Setting & Exact contributor fractions & \(R\) & \(w_U(t)\) \\
\midrule\noalign{}
\endhead
\bottomrule\noalign{}
\endlastfoot
G1 & \((1/3,1/3,1/3)\) & \(6.17284\) & \(0.201837\) \\
G4 & \((19/45,29/90,23/90)\) & \(63.6999\) & \(0.826856\) \\
G10 & \((3/5,3/10,1/10)\) & \(92.5924\) & \(0.99999875\) \\
\end{longtable}
}

Every setting has the same six positive channels. The quantitative
pattern changes. At G10,

\[
y=(0.05,0.05,0.15,0.15,0.30,0.30).
\]

Under the prescribed fractions and error model, the highest channels
strongly distinguish the target in the first position.

Equal-height observations do not require equal marginal genotype
likelihoods: population weighting of complementary co-contributor
configurations matters.

The target probability remains \(0.0108\), and

\[
R\leq\frac{1}{0.0108}\approx92.59.
\]

The final result approaches this common-model bound. The final
uniform-reference different-genotype comparison is approximately

\[
B=1.5974\times10^7.
\]

Its greater magnitude does not make it a better answer to the
different-person question. At fixed genotype-universe size, it is a
transformation of \(w_U(t)\).

\begin{figure}
\centering
\includegraphics[width=\linewidth]{figure-4.png}
\caption{The same alleles with changing quantitative patterns}
\end{figure}

\textbf{Figure 4. The same six allele channels can support different
quantitative discrimination in example G.} Three settings at one
synthetic locus are generated from \(5/6\), \(3/4\), and \(1/2\) in the
displayed fraction order. Generator truth is not treated as known
contributor information during evaluation. The target proposition fixes
\(5/6\) in the first position and population-weights the other
genotypes; the alternative population-weights all three positions.
Fractions are prescribed under both propositions. All channels have
positive signal, but their relative heights change. The target
probability remains \(0.0108\); \(w_U\) uses a uniform target-genotype
prior with population-weighted co-contributors and is not a personal
contribution probability. Gaussian channel errors have \(\sigma=0.035\).
Heights represent supplied signals; widths and shapes are schematic.
Fractions and statistics are rounded, and settings are not independent
loci.

\section{8. What each quantity
contributes}\label{what-each-quantity-contributes}

{\def\LTcaptype{none} % do not increment counter
\begin{longtable}[]{@{}
  >{\raggedright\arraybackslash}p{(\linewidth - 4\tabcolsep) * \real{0.3333}}
  >{\raggedright\arraybackslash}p{(\linewidth - 4\tabcolsep) * \real{0.3333}}
  >{\raggedright\arraybackslash}p{(\linewidth - 4\tabcolsep) * \real{0.3333}}@{}}
\toprule\noalign{}
\begin{minipage}[b]{\linewidth}\raggedright
Quantity
\end{minipage} & \begin{minipage}[b]{\linewidth}\raggedright
Question answered
\end{minipage} & \begin{minipage}[b]{\linewidth}\raggedright
What it does not supply by itself
\end{minipage} \\
\midrule\noalign{}
\endhead
\bottomrule\noalign{}
\endlastfoot
Population genotype probability & How frequent is the target genotype
under the declared distribution? & Contribution probability or
evidence-based genotype resolution \\
Population compatibility & What population mass belongs to a
rule-defined genotype set? & Identification of actual alternative
contributors \\
Genotype posterior & How is uncertainty distributed over genetic states
under a prior and model? & Automatic named-person attribution \\
Likelihood ratio & How do observations compare under stated
propositions? & Posterior odds without prior odds; for these
contributor-level comparisons, activity attribution or a legal
decision \\
\end{longtable}
}

The Gaussian model assigns positive density to every finite observation
under every genotype configuration. Small likelihoods in D and G are not
the strict model's structural impossibilities.

The examples therefore require different interpretive safeguards. A rare
target may receive no support when its position has no signal. A
strongly discriminating observation may approach the matching-genotype
bound. A large LR may coexist with diffuse genotype uncertainty. None of
those observations supplies the result for a different contributor
proposition.

\begin{center}\rule{0.5\linewidth}{0.5pt}\end{center}

\section{Part III --- Reporting the disputed
proposition}\label{part-iii-reporting-the-disputed-proposition}

\section{9. A bounded comparison of three
reports}\label{a-bounded-comparison-of-three-reports}

This part compares constructed reports under common evidence and
assumptions. It examines:

\begin{enumerate}
\def\labelenumi{\arabic{enumi}.}
\tightlist
\item
  Information sufficient for an informed reader to derive a conclusion.
\item
  Explicit warnings about the scope of reported comparisons.
\item
  Express reporting of the assessment relevant to the dispute.
\end{enumerate}

Actual reader understanding and legal adequacy are separate questions.
This is neither a comprehension experiment nor a survey of laboratory
practice.

The principal dispute is whether A and B form the complete two-person
configuration under the strict model.

\subsection{9.1 Common information}\label{common-information}

All versions receive:

\begin{itemize}
\tightlist
\item
  The observed allele sets and both reference genotypes.
\item
  The strict two-contributor observation rule.
\item
  The fixed-frequency independent population model.
\item
  The exclusion of A and B as unknown identities, not genotypes.
\item
  The selected-scenario status of the separate comparisons.
\item
  The independent-locus assumption for full-profile values.
\end{itemize}

All therefore contain the facts needed to derive the named-pair
incompatibility.

The quoted passages are read with this common information. They are not
standalone reporting templates.

\subsection{9.2 Version A --- Accurate baseline
report}\label{version-a-accurate-baseline-report}

\begin{quote}
Under the stated model, the observed thirteen-locus allele sets are
approximately \(6.69\times10^{13}\) times as likely if A and one unknown
person contributed as if two unknown people contributed.

Separately, the observations are approximately \(3.80\times10^{14}\)
times as likely if B and one unknown person contributed as if two
unknown people contributed.

Each result compares the likelihood of the observations under its stated
propositions. Neither number is a probability that A or B contributed,
or a probability of guilt.
\end{quote}

This constructed baseline has not been established as representative of
a particular laboratory's practice.

\subsection{9.3 Version B --- Baseline plus explicit
interpretation}\label{version-b-baseline-plus-explicit-interpretation}

Version B includes A and adds:

\begin{quote}
The two results concern different configurations. The unknown partner in
A's comparison need not be the partner in B's comparison. Neither
separate numerator proposes A and B together.

Consequently, the size of these separate results does not determine the
result of a comparison concerning A and B as the complete contributor
pair.

Another person may share a named person's genotype. These calculations
must therefore not be described as identifying a unique person merely by
distinguishing a genotype.

Even establishing DNA contribution would not, by itself, establish the
activity responsible for deposition or guilt.
\end{quote}

B adds interpretation without additional numerical results.

\subsection{9.4 Version C --- Clarification plus the joint
assessment}\label{version-c-clarification-plus-the-joint-assessment}

Version C includes B and adds:

\begin{quote}
We also evaluated A and B as the complete two-person contributor
configuration, against the same two-unknown alternative.

At each locus, A's genotype \(2/3\) requires a partner carrying \(1/4\).
B's genotype \(1/3\) requires a partner carrying \(2/4\).

Together, A and B supply alleles \(1\), \(2\), and \(3\), but not
observed allele \(4\).

The named pair therefore has zero likelihood under this strict model,
and its joint LR is \(0\).

This result concerns A and B as the complete two-person configuration.
It does not exclude their contribution somewhere within a larger mixture
or under a different observation model.
\end{quote}

C reports the joint assessment, its basis, and its limits.

\section{10. Reporting analysis and a bounded
recommendation}\label{reporting-analysis-and-a-bounded-recommendation}

\subsection{10.1 What C adds to B}\label{what-c-adds-to-b}

All three versions disclose the common information needed to derive the
joint incompatibility. None adds new biological observations.

Version A reports accurate selected comparisons but does not expressly
identify their joint limitation. Version B makes a substantive
improvement: it warns that the separate numerators evaluate different
configurations and do not determine the result for the named pair. B
must not be grouped with a presentation that invites an unqualified
joint-contribution inference.

C goes further. It expressly states that the joint configuration was
assessed, reports the incompatibility, explains the missing allele, and
identifies the model-specific scope. Its increment over B is the joint
conclusion and its basis, not the first warning against combining the
separate results.

An explicit assessment identifies the assessor's evaluation rather than
requiring the recipient to reconstruct both the result and whether the
disputed proposition was evaluated. That is a transparency rationale. It
does not establish that C improves comprehension or that A or B is
legally inadequate.

\subsection{10.2 Relevance and
factorization}\label{relevance-and-factorization}

The relevant assessment depends on the dispute. If A's contribution is
accepted and B's disputed, comparing A and B with A and an unknown
addresses a different question from comparing A and B with two unknowns.
This does not justify assuming A's contribution when it remains
disputed.

The separate LRs in Section 3 cannot simply be multiplied to obtain the
named-pair LR. Their numerators describe different configurations and do
not provide the conditional likelihood factors required for that
calculation.

This does not deny valid factorization among appropriately matched
compound and conditional comparisons, as discussed by Wivell et
al.~(2023, Section 2.3 and Appendix B). It rejects multiplication of
these particular selected results as independent factors for the named
pair.

Nor is the joint zero universally ``defense evidence.'' Its legal
significance must be tied to the proposition challenged; it does not
identify which individual contributed or establish either person's
exoneration.

\subsection{10.3 Existing reporting
recommendations}\label{existing-reporting-recommendations}

The distinction between evaluation and communication is analytically
useful but is not itself a novelty claim. Wivell et al.~(2023, Section
5) expressly state that when candidates cannot jointly be donors, or the
compound LR strongly disfavors their joint contribution, ``it is still
necessary to state this explicitly.''

Slooten (2022, Introduction) advocates keeping the relevant hypothesis
likelihoods central and reporting their table, while allowing individual
LRs as well. That approach can preserve information obscured by isolated
scalar comparisons.

The ASB 041 ballot draft addresses limitations of selected propositions
and their possible divergence from parties' positions, joint assessment,
and reasonable reassessment requests. Clause 4.14 requires documenting
why a subset of evaluative LRs is selected for reporting and noting the
existence of additional LRs in the case file; it does not require every
numerical result to appear in the report. Its Foreword discusses
avoiding unintended attribution of propositions to prosecution or
defense.

These are direct precedents. Acknowledging them does not require
endorsing every proposed conditioning practice or assuming a disputed
person contributed. The contribution here is the linked analysis of
different interpretive problems and the distinct responses they require.

\subsection{10.4 Bounded transparency
workflow}\label{bounded-transparency-workflow}

For this recommendation, an assessment is material when it addresses a
proposition actually advanced or contested, or supplies a limitation
needed to avoid overstating the meaning attributed to a reported result.
The proposition or attributed meaning must be identified. This is a
working criterion, not a legal materiality test.

The recommendation does not require an analyst to anticipate every
possible defense or require a defendant to disclose a theory.
Responsibility is bounded by the question communicated, the work
performed, the method's capabilities, and relevant information
subsequently received.

Issuing a report, giving testimony, and advocating from another person's
assessment are distinct communicative acts. Under this recommendation,
responsibility attaches to the assessment and meaning a person
communicates in context. A later speaker's overstatement is not
automatically an analytical or reporting error by the original assessor.
When a broader proposition is subsequently put to the assessor, however,
the assessor should identify whether the existing evaluation addresses
it and communicate any established limitation material to that question.

For example, if separate results are used to support A and B as the
complete named pair, their established joint incompatibility directly
limits that attributed meaning. If only a selected
individual-with-unknown comparison is at issue and no broader meaning is
identifiable, the mere existence of a known joint result does not
automatically make it material. The recommendation distinguishes
communicating an established limitation from undertaking a new
evaluation.

{\def\LTcaptype{none} % do not increment counter
\begin{longtable}[]{@{}
  >{\raggedright\arraybackslash}p{(\linewidth - 2\tabcolsep) * \real{0.5000}}
  >{\raggedright\arraybackslash}p{(\linewidth - 2\tabcolsep) * \real{0.5000}}@{}}
\toprule\noalign{}
\begin{minipage}[b]{\linewidth}\raggedright
Situation
\end{minipage} & \begin{minipage}[b]{\linewidth}\raggedright
Recommended reporting response
\end{minipage} \\
\midrule\noalign{}
\endhead
\bottomrule\noalign{}
\endlastfoot
A material limitation is already established & State the limitation, its
basis, and model-specific scope. Do not leave an incompatible broader
interpretation unqualified. \\
A relevant comparison has not been evaluated & Identify non-evaluation
where needed to prevent the result from being understood as answering
that question. Non-evaluation is not an adverse result. \\
A reasonable later request or changed proposition arises & Provide a
route for reassessment. Identify changed assumptions or comparisons and
explain whether earlier conclusions retain their scope. \\
The available method cannot validly assess the comparison & State the
limitation. Do not substitute a computable but materially different
comparison without identifying the difference. \\
A speculative configuration has not been made material & No general
obligation to calculate it follows. Reconsider that status if later
information makes it relevant. \\
\end{longtable}
}

The practical sequence is to record the question, identify the
propositions assessed, state established material limitations, identify
relevant non-evaluation or methodological inability, and provide a route
for reasonable reassessment.

A numerical LR is not always necessary. In the strict example,
explaining that the complete named pair cannot supply an observed allele
communicates the essential incompatibility. If a numerical comparison is
given, its alternative and model must be clear.

The appropriate detail may differ among a technical report,
lawyer-facing explanation, testimony, and a jury-facing exhibit. The
figures here are scholarly illustrations, not validated exhibits for
those other settings.

\section{11. The separate role of genotype
uncertainty}\label{the-separate-role-of-genotype-uncertainty}

The genotype-posterior discussion addresses claims about genotype
resolution. It is not an alternative route to a probability of
innocence.

Two potential supplements remain distinct:

\begin{itemize}
\tightlist
\item
  An express assessment of the complete named pair.
\item
  A description of genotype uncertainty under a declared prior and
  model.
\end{itemize}

A--C examines the first. It does not demonstrate the practical or legal
value of the second.

A genotype posterior may expose an unsupported assertion that the
mixture resolves a particular genotype. It does not follow that every
large LR requires a posterior display. Alternative-genotype likelihoods
and the matching-genotype branch may already be included in the LR
denominator, even when their distribution is not displayed.

Let \(G_j\) denote the genotype in position \(j\), and let POI mean
person of interest. The alternative \(G_j\ne t\) does not exhaust
non-contribution by the POI: another person can carry \(t\). If
contribution anywhere is disputed, uncertainty over positions also
matters.

Similarly, a relative's possession of some alleles does not establish a
matching profile or a compatible mixture explanation. Under ordinary
inheritance without mutation, parents collectively carry alleles
transmitted to a child, but neither necessarily has the child's profile,
and a parental mixture generally includes additional alleles. Any such
contributor explanation requires its own evaluation. Appendix B.4
provides a limited transmission illustration.

These distinctions identify when uncertainty may matter to an attributed
meaning. They do not establish a general requirement to eliminate every
alternative genotype before DNA evidence supports contribution. Weak
discrimination may be uninformative rather than strongly exclusionary.

\section{12. McDaniel: legal context and an interpretive
hypothesis}\label{mcdaniel-legal-context-and-an-interpretive-hypothesis}

In McDaniel v. Brown, the Supreme Court distinguished random match
probability from the probability that the defendant was not the DNA
source, and distinguished source attribution from guilt, 558 U.S. 120,
128 (2010). It also emphasized the importance of fair and reliable
presentation of persuasive DNA evidence, id. at 136. These statements
provide context for careful probabilistic communication; they do not
prescribe this article's reporting format.

The procedural setting matters. The Court required consideration of all
admitted trial evidence in the sufficiency analysis, id. at 131. It did
not resolve the experts' relative credibility and explained that the
claim failed even assuming Mueller's estimate, id. at 132 and n.5. It
treated the separately advanced DNA due-process claim as forfeited and
reversed and remanded, with ineffective-assistance issues remaining, id.
at 134--136.

Justice Thomas, joined by Justice Scalia, objected to the extended
discussion of the post-trial Mueller report because it had no place in
the relevant sufficiency inquiry, id. at 137--138. He declined to join
the discussion he characterized as dicta concerning that report's effect
on the constitutional analysis. That objection should not be described
as a general rejection of accurate probabilistic communication.

An RMP is a probability, whereas \(H_d\) is a proposition. In an
appropriately specified ideal resolved-match model, the RMP may equal
\(\Pr(E\mid H_d,I)\). It does not thereby equal \(\Pr(H_d\mid E,I)\).
Nor does its use establish that the selected scientific alternative
represents a defendant's adopted explanation or exhausts relevant
non-contribution scenarios. A logical complement also does not supply
every distribution required to calculate its evidence likelihood.

We do not attribute to McDaniel a holding that an RMP-based likelihood
ratio is improper, or that the Court prescribed the scope of a defense
proposition. We advance a narrower interpretive hypothesis: its concern
with assigning unsupported meaning to DNA statistics may also bear on
presenting a restricted scientific comparison as though it evaluates a
materially different defense explanation.

The proposed connection has a reasoned basis. Probability transposition
and overstatement of evaluative scope can both assign a statistic a
meaning its calculation does not establish. They are nevertheless
different errors, and a decision addressing one does not automatically
establish a legal principle governing the other. General language about
fair presentation is not sufficient by itself to resolve that question.

If supported, the extension could matter even where the reported
statistic is correctly calculated: the concern would be the broader
assessment attributed to it. Neither that legal extension nor its
application to an actual report is established here. The proposed
connection does not establish an admissibility rule, a constitutional
violation, or a basis for relief in any particular proceeding; those
conclusions would require additional authority and case-specific
analysis. The mathematical results and independent transparency
rationale do not depend on acceptance of the hypothesis.

\section{13. Synthesis and conclusion}\label{synthesis-and-conclusion}

The integrated analysis separates problems that should not be addressed
as though they were interchangeable.

{\def\LTcaptype{none} % do not increment counter
\begin{longtable}[]{@{}
  >{\raggedright\arraybackslash}p{(\linewidth - 2\tabcolsep) * \real{0.5000}}
  >{\raggedright\arraybackslash}p{(\linewidth - 2\tabcolsep) * \real{0.5000}}@{}}
\toprule\noalign{}
\begin{minipage}[b]{\linewidth}\raggedright
Attributed meaning that requires examination
\end{minipage} & \begin{minipage}[b]{\linewidth}\raggedright
Relevant response
\end{minipage} \\
\midrule\noalign{}
\endhead
\bottomrule\noalign{}
\endlastfoot
Separate support is treated as joint support & Evaluate the joint
configuration or identify its non-evaluation. \\
A selected comparison is treated as an exhaustive individual-attribution
assessment & Identify the additional scenarios and distributions
required for a composite comparison. \\
A large LR is treated as genotype resolution & Examine genotype
uncertainty under the declared prior and model. \\
Derivable information is treated as an expressly reported assessment &
State the material conclusion and its scope. \\
A strict-model impossibility is transferred to another observation
mechanism & Reevaluate under that mechanism rather than importing the
zero. \\
\end{longtable}
}

This is the practical purpose of the integrated mathematical treatment.
The strict construction demonstrates a proposition mismatch and, through
the matched posterior identity, a separate resolution distinction. D and
G show why rarity alone does not determine discrimination. A--C isolates
an explicit-reporting difference without changing biological
information.

The six-face organization is useful insofar as it tracks dependencies.
Changing the contributor count can change likelihoods and exclusion
language. Changing the prior can change a genotype posterior without
changing conditional likelihoods. Changing the dispute can change which
comparison matters. Changing population or relationship assumptions
requires reconsidering the affected distributions rather than applying
an unexplained final-number adjustment.

The analysis does not establish how frequently these issues arise in
practice, whether a particular report creates an unsupported inference,
or whether additional explanation improves understanding. Actual
reporting records, operational validation, implementation documentation,
participant research, or case-specific legal analysis would be required
for those different claims. No probabilistic-genotyping export or
population-conditioning implementation was audited here.

Correct calculation does not, by itself, establish that an assessment
addresses the disputed proposition or supports the meaning attributed to
it. That question concerns the calculation's propositions and
assumptions as well as its communicated scope. Equally, a mathematical
illustration alone does not establish impaired fairness. Existing
practices may already handle the relevant distinctions, and a report
need not contain every computable statistic to communicate adequately.

The linked results make the distinction concrete. Strong selected
separate support does not establish complete-pair contribution; large
relative support does not establish genotype resolution; and disclosing
facts sufficient to derive a conclusion is not identical to expressly
reporting its assessment. These are different limitations, not
cumulative proof of one operational defect. The proposed response is to
identify the material question, determine whether the assessment
performed addresses it, and communicate the conclusion and its limits.

\section{14. Methods and numerical
reproducibility}\label{methods-and-numerical-reproducibility}

\subsection{14.1 Strict enumeration}\label{strict-enumeration}

The strict observation likelihood is one for configurations whose
allele-set union equals the observed complete set and zero otherwise.
Unknown genotypes are independent Hardy--Weinberg draws from the fixed
frequencies. Full-profile products use stipulated locus independence.

The accompanying implementation enumerates fifteen unordered genotypes
and all 225 ordered contributor pairs using exact rational arithmetic.
It checks the denominator, separate numerators, named-pair
incompatibility, designated-position posterior, profile count, and
matched posterior--frequency identities.

Reference independence is a separate stipulation, not a consequence of
Hardy--Weinberg proportions alone.

\subsection{14.2 Quantitative likelihood}\label{quantitative-likelihood}

Each allele copy contributes half its contributor fraction to the
corresponding channel:

\[
\mu_a=\sum_j\frac{\phi_j}{2}n_a(g_j),
\]

where \(n_a(g_j)\) counts allele copies.

Observations are deterministic signals constructed from the generating
genotypes and prescribed fractions. Evaluation uses

\[
K(y\mid\mathbf g,\boldsymbol\phi)
\propto
\exp\left[
-\frac{\sum_a(y_a-\mu_a)^2}{2\sigma^2}
\right],
\qquad
\sigma=0.035.
\]

The omitted Gaussian factor is common to the configurations and
comparisons in each setting.

For D, with known major genotype \(m=1/2\),

\[
L_D(g)=K(y\mid m,g;\boldsymbol\phi).
\]

For G,

\[
L_G(g)=
\sum_{h,k\in\mathcal G}
\pi(h)\pi(k)
K(y\mid g,h,k;\boldsymbol\phi).
\]

The quantities \(R\), \(w_U(t)\), \(w_\pi(t)\), and \(B\) are then
calculated from the common \(L_D\) or \(L_G\) using Section 4's
definitions. All twenty-one target-position genotypes are included; G
enumerates all ordered co-contributor pairs.

This model has no stutter, censoring, degradation, laboratory dropout,
or drop-in mechanism. It permits negative observations and has positive
density at every finite signal vector. It is a mathematical toy model,
not validated casework software.

\subsection{14.3 Exact inputs}\label{exact-inputs}

All population frequencies and \(\sigma\) stated above are treated as
exact decimals before floating-point likelihood evaluation.

Signal arrays use allele-channel order \(1\) through \(6\).

{\def\LTcaptype{none} % do not increment counter
\begin{longtable}[]{@{}
  >{\raggedright\arraybackslash}p{(\linewidth - 4\tabcolsep) * \real{0.3333}}
  >{\raggedright\arraybackslash}p{(\linewidth - 4\tabcolsep) * \real{0.3333}}
  >{\raggedright\arraybackslash}p{(\linewidth - 4\tabcolsep) * \real{0.3333}}@{}}
\toprule\noalign{}
\begin{minipage}[b]{\linewidth}\raggedright
Setting
\end{minipage} & \begin{minipage}[b]{\linewidth}\raggedright
Exact fraction vector
\end{minipage} & \begin{minipage}[b]{\linewidth}\raggedright
Exact signal array
\end{minipage} \\
\midrule\noalign{}
\endhead
\bottomrule\noalign{}
\endlastfoot
D1 & \((1,0)\) & \((1/2,1/2,0,0,0,0)\) \\
D6 & \((8/9,1/9)\) & \((4/9,4/9,0,0,1/18,1/18)\) \\
D10 & \((4/5,1/5)\) & \((2/5,2/5,0,0,1/10,1/10)\) \\
G1 & \((1/3,1/3,1/3)\) & \((1/6,1/6,1/6,1/6,1/6,1/6)\) \\
G4 & \((19/45,29/90,23/90)\) &
\((23/180,23/180,29/180,29/180,19/90,19/90)\) \\
G10 & \((3/5,3/10,1/10)\) & \((1/20,1/20,3/20,3/20,3/10,3/10)\) \\
\end{longtable}
}

D's fraction order is known major \(1/2\), then target-position genotype
\(5/6\). G's generating order is \(5/6\), \(3/4\), \(1/2\).

\subsection{14.4 Implementation and numerical
results}\label{implementation-and-numerical-results}

The implementation uses only the Python standard library. It uses exact
rational arithmetic for the strict construction and explicit rational
inputs for the quantitative examples. Gaussian calculations use
floating-point log likelihoods and log-sum-exp marginalization.

{\def\LTcaptype{none} % do not increment counter
\begin{longtable}[]{@{}
  >{\raggedright\arraybackslash}p{(\linewidth - 2\tabcolsep) * \real{0.4286}}
  >{\raggedleft\arraybackslash}p{(\linewidth - 2\tabcolsep) * \real{0.5714}}@{}}
\toprule\noalign{}
\begin{minipage}[b]{\linewidth}\raggedright
Quantity
\end{minipage} & \begin{minipage}[b]{\linewidth}\raggedleft
Numerical result
\end{minipage} \\
\midrule\noalign{}
\endhead
\bottomrule\noalign{}
\endlastfoot
Strict denominator & \(0.0012096\) \\
Strict A numerator & \(0.014\) \\
Strict B numerator & \(0.016\) \\
Strict joint LR & \(0\) \\
Strict A thirteen-locus LR & \(6.68838628684560\times10^{13}\) \\
Strict B thirteen-locus LR & \(3.79504262792686\times10^{14}\) \\
Strict target-position profile posterior &
\(7.65656096662972\times10^{-11}\) \\
Compatible position profiles & \(13{,}060{,}694{,}016\) \\
\end{longtable}
}

{\def\LTcaptype{none} % do not increment counter
\begin{longtable}[]{@{}lrr@{}}
\toprule\noalign{}
Setting & \(R\) & \(w_U(t)\) \\
\midrule\noalign{}
\endhead
\bottomrule\noalign{}
\endlastfoot
D1 & \(1\) & \(0.0476190476190476\) \\
D6 & \(27.9507017179653\) & \(0.541728524905921\) \\
D10 & \(91.9452838579187\) & \(0.997158306143090\) \\
G1 & \(6.17283949624454\) & \(0.201837236145844\) \\
G4 & \(63.6999391814617\) & \(0.826855819812353\) \\
G10 & \(92.5923708630398\) & \(0.999998747977069\) \\
\end{longtable}
}

At D6, the population-prior target posterior was \(0.301867578554025\).
The final G value of \(B\) was \(15{,}974{,}128.3188332\).

The recorded numerical results were obtained using CPython 3.11.14. All
132 programmed checks passed: 22 for the strict construction and
population filters, 78 for the six D/G settings, and 32 for supporting
examples. These are related assertions, not 132 independent validations.

The strict calculations use exact comparisons where applicable.
Main-table D/G values are checked against their displayed rounding, with
a small floating-point allowance. Higher-precision comparisons and
numerical identities use explicit tolerances recorded in the script. The
tested values agree with the displayed results.

The complete code, running instructions, and numerical-results record
are included in the supplementary material.

\subsection{14.5 Scope of the numerical
checks}\label{scope-of-the-numerical-checks}

Supporting checks combine finite enumeration with direct formula
evaluation. In particular, the source-posterior and occupancy examples,
fifteen-locus population-filter products, and conditional-relatedness
arithmetic are formula checks, not independent validation of their
modeling assumptions. The programmed checks do not certify every general
derivation or every possible application.

The quantitative calculations cover the six D/G settings specified in
Section 14.3. They do not validate commercial software, actual
laboratory observations, reader comprehension, or legal claims.

The reporting analysis compares constructed texts under common facts. It
contains no participant observations.

Figure widths and shapes are illustrative; the statistical model uses
channel heights, not drawn peak geometry.

\begin{center}\rule{0.5\linewidth}{0.5pt}\end{center}

\section{Appendix A. Genotype posteriors and uniform-reference
comparisons}\label{appendix-a.-genotype-posteriors-and-uniform-reference-comparisons}

\subsection{A.1 Posterior recovery and marginal
priors}\label{a.1-posterior-recovery-and-marginal-priors}

For a complete posterior under known positive prior \(q(g)\),

\[
w(g)=\frac{q(g)L(g)}{\sum_hq(h)L(h)},
\qquad
L(g)\propto\frac{w(g)}{q(g)}.
\]

Under a known positive prior, uncorrected posterior weights are
proportional to likelihoods when the prior is constant over
positive-likelihood states. Uniformity over zero-likelihood states is
unnecessary.

State definitions, conditioning, and nuisance treatment must agree.
Unknown omitted tails or incompatible aggregation prevent unrestricted
recovery. Probabilistic-genotyping systems use differing approaches,
making model-specific interpretation essential (Gill et al., 2021). No
actual software export was examined here.

A flat joint distribution need not induce a flat marginal. With two
labelled contributors whose genotypes are AA, AB, and BB, assign equal
mass to the seven ordered configurations AA/AB, AA/BB, AB/AA, AB/AB,
AB/BB, BB/AA, and BB/AB. The first contributor's marginal probabilities
are \(2/7\), \(3/7\), and \(2/7\).

With a constant subsequent likelihood, those remain the posteriors.
Treating them as uniform-prior likelihoods falsely favors AB. Dividing
by the actual prior recovers equal likelihoods.

If configurations were selected using the evidence being evaluated,
their distribution is conditional on that selection and cannot be reused
as an unconditional prior for the same evidence.

\subsection{A.2 Uniform-reference comparison and a
counterexample}\label{a.2-uniform-reference-comparison-and-a-counterexample}

For \(M>1\) atomic states,

\[
B=\frac{(M-1)L(t)}{\sum_{g\ne t}L(g)}
=\frac{(M-1)w_U(t)}{1-w_U(t)}.
\]

Its inverse on the finite interior domain is

\[
w_U(t)=\frac{B}{B+M-1}.
\]

This is an application of the posterior-odds identity, not additional
evidence (Kass and Raftery, 1995). At fixed \(M\), rankings and
receiver-operating-characteristic curves remain unchanged under
corresponding threshold transformation. The same numerical threshold on
different scales need not give the same decisions.

In general, \(B\) and \(f(t)\) do not determine \(R\). Take three
genotypes with population probabilities \(0.1\), \(0.8\), and \(0.1\),
the first being the target. Likelihoods \(0.6\), \(0.3\), and \(0.1\)
give \(B=3\) and \(R\approx1.935484\). Swapping the last two likelihoods
leaves \(B\) and \(f(t)\) unchanged but gives \(R\approx3.529412\).

The population weighting, not merely the sum of alternative likelihoods,
matters.

\subsection{A.3 Universe definition and prior
dilution}\label{a.3-universe-definition-and-prior-dilution}

Adding zero-likelihood states while holding existing likelihoods fixed
increases \(B\) by spreading the uniform alternative's mass more thinly.
This is prior sensitivity, not new evidence. Its relationship to
diffuse-prior discussions should not be mistaken for independent data
support (Bartlett, 1957).

For twenty toy loci with twenty-five alleles per locus, each locus has
\(325\) unordered genotypes. The product of twenty factors of \(324\) is
approximately

\[
324^{20}\approx1.6251753\times10^{50}.
\]

This is the multiplier in a product of locus-specific comparisons
against genotypes differing at every locus, subject to factorization. It
is not the multiplier for a whole-profile alternative excluding only the
target profile. Equal likelihoods cancel the multiplier exactly.

A physical allele ladder is not automatically a complete biological
genotype universe.

\subsection{A.4 Structural support, truncation, and
precision}\label{a.4-structural-support-truncation-and-precision}

If \(k>1\) states have positive likelihood, including \(t\), then \(B\)
factors into \((M-1)/(k-1)\) times the ratio of the target likelihood to
the mean positive alternative likelihood. When every genotype has
positive likelihood, \(k=M\).

A threshold is not a structural zero. If the retained alternative
likelihood sum is \(A\) and the omitted sum is \(T\),

\[
B=\frac{(M-1)L(t)}{A+T}.
\]

For \(A>0\), this equals the truncated result multiplied by \(A/(A+T)\).
If \(A=0<T\), calculate directly rather than multiplying an infinite
value by zero.

Aggregation preserves the comparison only when atomic prior masses and
likelihood contributions are retained. Treating displayed categories as
newly equally weighted states changes the model.

Positive target likelihood with zero alternative likelihood gives an
infinite \(B\) under an extended-real convention. If all likelihoods are
zero, the ratio is undefined.

A displayed posterior of \(1.0000\) establishes neither condition. With
nearest rounding of a complete posterior under a uniform target-genotype
prior and \(M=55\), it implies only the lower bound

\[
B\geq1{,}079{,}946.
\]

That bound does not apply to a posterior renormalized after truncation.

\subsection{A.5 Alternative-specific
expectations}\label{a.5-alternative-specific-expectations}

Consider fifty-five genotypes and fifty-five observations. Conditional
on a genotype, its corresponding observation has probability \(0.84\);
each other observation has probability \(0.16/54\). Fix the target
before observing the evidence.

The uniform-reference \(B\) is \(283.5\) on the corresponding
observation and approximately \(0.160475483\) otherwise. Under its
uniform different-genotype alternative, the first event has probability
\(0.002962963\), and the mean of \(B\) is one.

Under a different-person alternative allowing all fifty-five genotypes
uniformly, including the target, that event has probability \(1/55\),
and the mean of \(B\) is approximately \(5.312103\).

The corresponding person-level LR is \(46.2\) on the target observation
and approximately \(0.162962963\) otherwise, with mean one under that
alternative.

An expectation property under one alternative therefore does not
automatically transfer to another. This example does not establish
comprehensive calibration or operational performance.

\subsection{A.6 Multilocus alternatives and descriptive
summaries}\label{a.6-multilocus-alternatives-and-descriptive-summaries}

Multiplying locus-specific uniform-reference comparisons tests the
alternative that the genotype differs at every locus, provided
likelihoods and alternative weights factorize. It does not test
difference somewhere in the profile.

With two loci, two genotypes per locus, independent uniform priors, and
factorizing likelihoods, target posterior \(0.8\) at each locus gives
joint posterior \(0.64\). The every-locus alternative gives Bayes factor
\(16\); the uniform joint alternative excluding only the target profile
gives \(5.333333\).

Genetic independence does not establish marginal-likelihood
factorization when nuisance parameters are shared across loci.

Target weight, rank, top-three mass, entropy, and reciprocal
squared-weight effective counts can describe genotype uncertainty. They
remain prior- and model-dependent; aggregation can change them.
Multiplying posterior odds by an effective count creates no LR
interpretation by itself.

\section{Appendix B. Supporting population and nuisance
examples}\label{appendix-b.-supporting-population-and-nuisance-examples}

\subsection{B.1 Occupancy and posterior source
attribution}\label{b.1-occupancy-and-posterior-source-attribution}

Under \(K-1\) independent alternative persons, each with matching
probability \(f\), the probability that at least one matches is

\[
1-(1-f)^{K-1}.
\]

This is a pre-observation occupancy probability, not the probability
another person contributed.

For a separate idealized source model, suppose exactly one of \(K\)
exchangeable candidates is the source, the person of interest was
identified independently of the questioned DNA, and the other candidate
profiles remain unobserved independent draws. For an error-free resolved
match,

\[
\Pr(P\mid E,I)=\frac{1}{1+(K-1)f},
\]

where \(P\) denotes the person of interest being the source.

For \(f=10^{-6}\), the posterior is approximately \(0.999901\) for
\(100\) candidates, \(0.990100\) for \(10{,}000\), and \(0.50000025\)
for one million. With one million candidates, the pre-observation
probability of another matching carrier is approximately \(0.632120\).

The quantities differ because their events and conditioning differ. This
is not an instruction to determine prior odds by counting everyone who
might conceivably have been nearby.

The ideal resolved-match LR is \(1/f\). It incorporates match-frequency
information but does not supply posterior source probability without
prior odds.

\subsection{B.2 Binary-mixture nuisance
treatment}\label{b.2-binary-mixture-nuisance-treatment}

Consider observed alleles \(7\), \(8\), \(9\), and \(10\) with
frequencies \(0.07\), \(0.08\), \(0.09\), and \(0.10\), and target
\(8/9\).

Target probability is \(0.0144\), allele-set inclusion is \(0.1156\),
and strict compatibility is \(0.0862\). Under independent population
draws, the target-plus-unknown likelihood is \(0.014\), the two-unknown
likelihood is \(0.0012096\), and the LR is approximately \(11.574074\).

For a separate uniform-reference illustration, define ten alleles: the
four observed and six unobserved alleles collectively carrying mass
\(0.66\). They form fifty-five genotypes. The allocation among
unobserved alleles does not affect this strict likelihood because
configurations containing them are incompatible.

With a uniform target-genotype prior and population-weighted
co-contributor, the six positive likelihoods are the
complementary-genotype frequencies

\[
0.0180,\quad0.0160,\quad0.0144,\quad0.0140,\quad0.0126,\quad0.0112.
\]

They give \(B\approx10.470914\). Making the co-contributor uniform
instead equalizes the six likelihoods and gives \(B=10.8\).

These are different nuisance models, not interchangeable calculations of
one comparison.

\subsection{B.3 Multilocus population
filters}\label{b.3-multilocus-population-filters}

Repeating Section 5's frequencies at fifteen independent loci with
target \(2/3\) gives:

{\def\LTcaptype{none} % do not increment counter
\begin{longtable}[]{@{}
  >{\raggedright\arraybackslash}p{(\linewidth - 4\tabcolsep) * \real{0.2727}}
  >{\raggedleft\arraybackslash}p{(\linewidth - 4\tabcolsep) * \real{0.3636}}
  >{\raggedleft\arraybackslash}p{(\linewidth - 4\tabcolsep) * \real{0.3636}}@{}}
\toprule\noalign{}
\begin{minipage}[b]{\linewidth}\raggedright
Event
\end{minipage} & \begin{minipage}[b]{\linewidth}\raggedleft
One locus
\end{minipage} & \begin{minipage}[b]{\linewidth}\raggedleft
Fifteen loci
\end{minipage} \\
\midrule\noalign{}
\endhead
\bottomrule\noalign{}
\endlastfoot
Target genotype & \(0.0144\) & \(2.373763\times10^{-28}\) \\
Strict two-person compatibility & \(0.0862\) &
\(1.077976\times10^{-16}\) \\
Both alleles inside observed set & \(0.1156\) &
\(8.797667\times10^{-15}\) \\
\end{longtable}
}

The filters are nested but answer different questions. Multiplying a
probability by an explicitly modeled candidate count gives an expected
count, not identified alternative contributors or a source posterior.

This fifteen-locus illustration is separate from the thirteen-locus
strict construction.

\subsection{B.4 Conditional
relatedness}\label{b.4-conditional-relatedness}

Take three allele sets with population masses \(0.28\), \(0.29\), and
\(0.32\). Under independent allele draws, an unrelated person's
inside-set probability is

\[
(0.28\times0.29\times0.32)^2=0.000675168256.
\]

If paternal transmission is inside each set with probability one and the
other transmitted allele is an independent population draw, the child's
probability is

\[
0.28\times0.29\times0.32=0.025984.
\]

The unrelated probability is the square of the child probability under
these assumptions. This is not a universal relationship for relatives
and does not test parental contribution or identify a particular child.

\begin{center}\rule{0.5\linewidth}{0.5pt}\end{center}

\section{Data, code, and supplementary
materials}\label{data-code-and-supplementary-materials}

All examples use synthetic finite models. No confidential casework
records, laboratory exports, or participant data are used.

The model specifications and exact inputs are given in Section 14 and
the appendices. The \href{supplemental.md}{supplementary material}
contains the complete standard-library Python implementation,
instructions for running it, the numerical-results record, and a map
connecting the supporting material to the article. No external dataset
is required.

The accompanying vector figures are \texttt{figure-1.svg},
\texttt{figure-2.svg}, \texttt{figure-3.svg}, and \texttt{figure-4.svg}.
Their captions specify the represented quantities and model boundaries.
In Figures 2--4, \(f(t)\) denotes the target genotype population
probability, not a contribution posterior. Peak shapes and widths in
Figures 3--4 are schematic.

\textbf{Declarations}

\textbf{Monte Miller PhD}

\textbf{No funding declared.}

\textbf{Author asserts no competing interest.}

\textbf{Ethics and consent:} The study uses synthetic models and no
confidential casework or participant data.

\textbf{AI assistance:} AI tools assisted with analytical framing,
numerical implementation, drafting, figure programming,
literature-comparison drafting, and revision.

\section{References}\label{references}

AAFS Standards Board. (2021). \emph{Formulating Propositions for
Likelihood Ratios in Forensic DNA Interpretations}. ASB Standard 041,
First Edition 2021, ballot draft, file \texttt{041\_Std\_Ballot02.pdf}.
The cited document is a ballot draft; no binding or adoption status is
asserted.
https://www.aafs.org/sites/default/files/media/documents/041\_Std\_Ballot02.pdf

Bartlett, M. S. (1957). A comment on D. V. Lindley's statistical
paradox. \emph{Biometrika}, 44(3--4), 533--534.
https://doi.org/10.1093/biomet/44.3-4.533

Bieber, F. R., Buckleton, J. S., Budowle, B., Butler, J. M., and Coble,
M. D. (2016). Evaluation of forensic DNA mixture evidence: protocol for
evaluation, interpretation, and statistical calculations using the
combined probability of inclusion. \emph{BMC Genetics}, 17, 125.
https://doi.org/10.1186/s12863-016-0429-7

Duke, K., Cuenca, D., Myers, S., and Wallin, J. (2022). Compound and
Conditioned Likelihood Ratio Behavior within a Probabilistic Genotyping
Context. \emph{Genes}, 13(11), 2031.
https://doi.org/10.3390/genes13112031

Gill, P., Benschop, C., Buckleton, J., Bleka, Ø., and Taylor, D. (2021).
A review of probabilistic genotyping systems: EuroForMix, DNAStatistX
and STRmix. \emph{Genes}, 12(10), 1559.
https://doi.org/10.3390/genes12101559

Kass, R. E., and Raftery, A. E. (1995). Bayes factors. \emph{Journal of
the American Statistical Association}, 90(430), 773--795.
https://doi.org/10.1080/01621459.1995.10476572

McDaniel v. Brown, 558 U.S. 120 (2010).
https://www.govinfo.gov/content/pkg/USREPORTS-558/pdf/USREPORTS-558-120.pdf

National Research Council. (1996). \emph{The Evaluation of Forensic DNA
Evidence}. Washington, DC: National Academies Press.
https://doi.org/10.17226/5141

Slooten, K. (2022). The comparison of DNA mixture profiles with multiple
persons of interest. \emph{Forensic Science International: Genetics},
56, 102592. https://doi.org/10.1016/j.fsigen.2021.102592

Wivell, R., Kelly, H., Kokoszka, J., Daniels, J., Dickson, L.,
Buckleton, J., and Bright, J.-A. (2023). An Investigation into Compound
Likelihood Ratios for Forensic DNA Mixtures. \emph{Genes}, 14(3), 714.
https://doi.org/10.3390/genes14030714

\end{document}
